\problem{P01}{质量检测的价值}{逆向选择；信息设计}
卖方知道物品价值 $X\sim\operatorname{Unif}[0,1]$，买方估值为 $3X/2$。买方报一次价格 $p\ge0$，卖方当且仅当 $X\le p$ 时成交。买方可先支付费用 $c\ge0$ 做一次检测：他事先选择阈值 $t\in(0,1)$，随后仅得知 $X<t$ 或 $X\ge t$，再分别报价。求最优检测阈值、两种结果下的报价，并判断何时应购买检测。
\solution 无检测时，$0\le p\le1$ 的期望净交易收益为
\[
 \int_0^p(3x/2-p)\,\dd x=-p^2/4\le0;
\]
$p>1$ 时为 $3/4-p<0$。所以不交易可得最优收益零。
低质量结果下 $X$ 均匀于 $[0,t)$。报价 $p\le t$ 的收益非正；报 $p\ge t$ 时条件均值收益为 $3t/4-p<0$。故应不交易。
高质量结果下，只需考虑 $p\in[t,1]$，因为低于 $t$ 不成交、高于 $1$ 只会增加付款。将条件收益乘以该结果概率 $1-t$，得到无条件贡献
\[
 H_t(p)=\int_t^p(3x/2-p)\,\dd x
 =-\frac{p^2}{4}+tp-\frac{3t^2}{4}.
\]
这是关于 $p$ 的凹二次式，最优价格为 $p_t=\min\{2t,1\}$。检测带来的交易收益为
\[
 W(t)=\begin{cases}
 t^2/4,&0<t\le1/2,\\
 t-1/4-3t^2/4,&1/2\le t<1.
 \end{cases}
\]
第一段递增，第二段在 $t=2/3$ 处唯一最大。因此最优阈值为 $2/3$；低结果不买，高结果报价 $1$，毛收益为 $1/12$。扣除费用后，相对于不检测的改善为 $1/12-c$：费用低于 $1/12$ 时检测，高于时不检测，等于时两者无差异。
\insight 报价会筛选成交对象，检测又改变报价前的信息。先在每个信息分支优化报价，再优化检测规则，不能用无条件均值代替分支价值。
